\section{Simulating Magnetised Flows}
\label{sec:sims}

for computational convenience, we absorb the factor of 4pi into b

Gent: take over of a LSD from a SSD

\subsection{Conservative Methods}
\label{sec:sims:godunov}

\subsubsection{The Riemann Problem}
\label{sec:sims:riemann}

a whole new family of waves wake up, and they invite many new types of instabilities that lead to interesting dynamics 

Working in the ideal limit, and restricting our attention to 1D, we linearise \autoref{eqn:mhd:mass}-\ref{eqn:mhd:monopoles} about a uniform background state with background states \(\bar{\rho}\), \(\bar{\mVector{b}}\), and perturbed states \(\rho = \bar{\rho} + \delta{\rho}\), \(\mVector{u} = \delta{\mVector{u}}\) \(\mVector{b} = \bar{\mVector{b}} + \delta{\mVector{b}}\), which, in the reference frame of the fluid, \(\bar{\mVector{u}} = \mVector{0}\), yield the linearised system in matrix form
\begin{align}
    \mppFrac{(\delta{\mVector{S}_p})}{t}
        + \mTensor{J}_p \mppFrac{(\delta{\mVector{S}_p})}{x}
        = 0,
\end{align}
where the perturbed primitive states have been collected into a state vector \(\mVector{S}_p\), and the corresponding Jacobian matrix (capturing information about the couplings between quantities) is \(\mTensor{J}_p\),
\begin{align}
    \delta{\mVector{S}_p}
        \equiv \begin{bmatrix}
            \delta{\rho} \\
            \delta{u_x} \\
            \delta{u_y} \\
            \delta{u_z} \\
            \delta{b_y} \\
            \delta{b_z}
        \end{bmatrix}
    , \qquad
    \mTensor{J}_p
        \equiv \begin{bmatrix}
            0 & \bar{\rho} & 0 & 0 & 0 & 0 \\
            c_s^2 / \bar{\rho} & 0 & 0 & 0 & \bar{b}_y / (4\pi\bar{\rho}) & \bar{b}_z / (4\pi\bar{\rho}) \\
            0 & 0 & 0 & 0 & -\bar{b}_x / (4\pi\bar{\rho}) & 0 \\
            0 & 0 & 0 & 0 & 0 & -\bar{b}_x / (4\pi\bar{\rho}) \\
            0 & \bar{b}_y & -\bar{b}_x & 0 & 0 & 0 \\
            0 & \bar{b}_z & 0 & -\bar{b}_x & 0 & 0
        \end{bmatrix}.
\end{align}
The characteristic wavespeeds are given by the eigenvalues \(\lambda\) satisfying
\begin{align}
    \det(\mTensor{J}_p - \lambda\mTensor{I})
        = \rbrac{\lambda^2 - \frac{(\bar{b}_x)^2}{4\pi\bar{\rho}}} \rbrac{
            \lambda^4
            - \rbrac{c_s^2 + (\bar{b}_x)^2} \lambda^2
            + c_s^2\frac{(\bar{b}_x)^2}{4\pi\bar{\rho}}
        }
        = 0
        , \label{eqn:mhd-system:characteristics}
\end{align}
where
\begin{align}
    \rbrac{u_A}^2
        \equiv \frac{b^2}{4\pi\rho}.
\end{align}
is the Alfv\'en speed, which measures the speed at which transverse disturbances supported by magnetic tension, and \((\bar{u}_A)^2 \equiv (\bar{b})^2 / (4\pi\bar{\rho}),\) is its value on the background state.

\autoref{eqn:mhd-system:characteristics} naturally factorises into a quadratic and quartic polynomial in \(\lambda\), which reveals all six characteristic wavespeeds. First, Alfv\'en waves can propagate in both directions along magnetic fields
\begin{align}
    \lambda_A^{(\pm)}
        = \pm \frac{\bar{b}_x}{\sqrt{4\pi\bar{\rho}}},
\end{align}
and then there are fast and slow magnetosonic waves,
\begin{align}
    \lambda_f^{(\pm)}
        &= \pm \frac{1}{2}\rbrac{
            c_s^2
            + (\bar{b}_x)^2
            + \sqrt{
                \rbrac{c_s^2 + (\bar{b}_x)^2}^2
                - 4c_s^2\frac{(\bar{b}_x)^2}{4\pi\bar{\rho}}
            }
        }
    , \\
    \lambda_s^{(\pm)}
        &= \pm \frac{1}{2}\rbrac{
            c_s^2
            + (\bar{b}_x)^2
            - \sqrt{
                \rbrac{c_s^2 + (\bar{b}_x)^2}^2
                - 4c_s^2\frac{(\bar{b}_x)^2}{4\pi\bar{\rho}}
            }
        }
        ,
\end{align}
respectively. Note, however, that the entropy mode, which would be present in adiabatic MHD, is absent here under the isothermal closure adopted in \autoref{sec:modelling:mhd}. Therefore we only have six of the seven modes,
\begin{align}
    \lambda_f^{(-)} \le \lambda_A^{(-)} \le \lambda_s^{(-)}
    \le \lambda_s^{(+)} \le \lambda_A^{(+)} \le \lambda_f^{(+)}.
\end{align}

left and right leaning wrt (to the entropy mode and) co-moving frame of the fluid


\subsubsection{Maintaining Monopole-free Fields}
\label{sec:sims:ct}

Lorentz forces, for Ampere's law without displacement current
\begin{align}
    \frac{1}{c} \mVector{j}\times\mVector{b}
        &= \frac{1}{4\pi} (\mCurl{\mVector{b}})\times\mVector{b} \\
        &= \frac{1}{4\pi} \rbrac{
                \underbrace{
                    b^2 (\mVectorUnit{b}\cdot\nabla)\mVectorUnit{b}
                }_\text{tension}
                - \overbrace{
                    \frac{1}{2} \nabla_\perp b^2
                }^\text{magnetic pressure}
                - \underbrace{
                    (\nabla\cdot\mVector{b})\mVector{b}
                }_\text{\shortstack{spurious field-aligned\\ acceleration}}
            }
        ,
\end{align}

everything is approximated (even things you don't want to be), expect for what can't be.

powell's 8-wave scheme (1999): resolve extra wave, then use source term to advect magnetic monopoles (8th wave) away; can give the wrong shock jump conditions

brackbil and bams (1980), very expensive elliptic PDF, smooth discontinuities in b, cleaning

dender 2002, lagrangian multiplier

\subsection{Chasing Asymptotic Flow Regimes}

Resolution challenges

tearing instabilities
